% Generated draft: requires mathematical and source review.
\documentclass[11pt]{article}
\usepackage[margin=25mm]{geometry}
\usepackage{fontspec}
\setmainfont{Latin Modern Roman}
\usepackage{amsmath,amssymb,mathtools,unicode-math}
\setmathfont{Latin Modern Math}
\usepackage{enumitem,needspace,xurl,hyperref}
\setlength{\parindent}{0pt}
\setlength{\parskip}{5pt}
\setlength{\emergencystretch}{3em}
\newcommand{\N}{\mathbb{N}}
\newcommand{\Z}{\mathbb{Z}}
\newcommand{\Q}{\mathbb{Q}}
\newcommand{\R}{\mathbb{R}}
\newcommand{\C}{\mathbb{C}}
\newcommand{\F}{\mathbb{F}}
\newcommand{\A}{\mathbb{A}}
\newcommand{\PP}{\mathbb{P}}
\newcommand{\E}{\mathbb{E}}
\newcommand{\eps}{\varepsilon}
\newcommand{\dd}{\,\mathrm{d}}
\begin{document}
\section*{UnsolvedMath ID: 30003123}
\textbf{Source problem code:} catalogue entry
\subsection{Definitions and mathematical statement}
The supplied record asks the following exponential degree-growth question. Some conventions needed for a fully precise formulation are not specified in the record and are identified below.

For a candidate formulation, interpret a curve over \(\Q\) as a smooth, projective, geometrically integral curve \(X\), of genus \(g\). Let \(t\in\Q(X)\) be nonconstant, defining a morphism
\[
 t:X\longrightarrow\PP^1_{\Q},\qquad
 \nu=[\Q(X):\Q(t)]\geq2.
\]
These regularity and geometric integrality conventions require confirmation against the original setup, which is referenced but not reproduced in the original statement.

Fix an algebraic closure \(\overline{\Q}\). For each integer \(n\geq1\), choose a geometric point \(P_n\in X(\overline{\Q})\) with \(t(P_n)=n\). Write \(\Q(P_n)\subseteq\overline{\Q}\) for its field of definition and
\[
 K_N=\Q(P_1)\cdots\Q(P_N)
     =\Q(P_1,\ldots,P_N)
\]
for the compositum. A finite critical value of \(t\) means a point \(b\in\overline{\Q}\subset\PP^1(\overline{\Q})\) such that some point above \(b\) has ramification index greater than \(1\).

Assume at least one of the following source conditions holds:
\begin{enumerate}
\item \(t\) has a finite critical value in \(\overline{\Q}\setminus\Q\).
\item The geometric function-field extension \(\overline{\Q}(X)/\overline{\Q}(t)\) is ``nonabelian,'' in the terminology of the supplied statement.
\end{enumerate}
The second condition requires clarification: no Galois hypothesis or definition of ``nonabelian'' is supplied. In particular, the record does not specify whether it refers to a Galois extension with nonabelian Galois group, to the Galois group of the normal closure, or to another convention.

The question is whether there exists a positive constant \(c(g,\nu)\), depending only on \(g\) and \(\nu\), such that
\[
 [K_N:\Q]\geq\exp\!\bigl(c(g,\nu)N\bigr)
 \qquad\text{for all sufficiently large integers }N.
\]
The natural universal reading is that the same \(c(g,\nu)\) works for every admissible pair \((X,t)\) and every sequence of choices \((P_n)_{n\geq1}\), with a threshold \(N_0\) that may depend on \((X,t)\) and the sequence. The source does not explicitly settle the quantifier over point choices or the allowed dependence of \(N_0\); this reading therefore remains provisional. No integers corresponding to ramified fibers are excluded in the supplied statement.
\subsection{Short English statement}
Under either of two stated conditions on a map \(t:X\to\PP^1\), must the field generated by one point above each of \(1,\ldots,N\) have degree at least \(\exp(c(g,\nu)N)\)? The meaning of the nonabelian condition, the curve conventions, and the quantifiers over point choices need confirmation. This is a statement-only draft requiring human mathematical review.
\subsection{Sources}
\begin{itemize}
\item Local registry: \texttt{data/unsolvedmath/problems.json}, supplied record with permanent UnsolvedMath ID \(30003123\), title \emph{Exponential Degree Growth in Composita of Fibers}, exact source problem code \texttt{catalogue entry}; collection: Oberwolfach Reports Open Problems. The record gives proposed year \(2016\) and identifies no individual proposer.
\item Supplied report citation: \emph{Diophantische Approximationen} (2017), Oberwolfach Reports, DOI \texttt{10.4171/owr/2016/21}, pp.\ 1099--1187. \url{https://doi.org/10.4171/owr/2016/21}. The supplied background describes Oberwolfach Reports as documenting mathematical workshops and their open-problem discussions.
\item Supplied literature source, ``OWR report'': \url{https://ems.press/content/serial-article-files/46625?nt=1}.
\item Supplied literature source, ``Dvornicich--Zannier'': \url{https://www.numdam.org/item/ASNSP_1994_4_21_3_421_0/}.
\item Supplied literature source, ``Bilu--Luca'': \url{https://arxiv.org/abs/1606.09164}.
\item Literature assessment dated exactly \texttt{2026-08-21}: the supplied curation labels the problem partially solved and states that the strong diversity conjecture is proved for degree-two and degree-three critical values and other special cases, while the general Schinzel formulation remains open. This is supplied dated curation derived from the OWR material, not an independently verified assessment. The cited claims and time-sensitive status require independent verification.
\item The registry labels the statement \texttt{corrected\_verified} and reports that mathematical content was preserved while notation, typography, or omitted definitions were repaired. These are supplied editorial assertions, not mathematical verification performed for this draft.
\end{itemize}
\end{document}
